% Part: first-order-logic
% Chapter: tableaux
% Section: identity

\documentclass[../../../include/open-logic-section]{subfiles}

\begin{document}

\olfileid[fr]{fol}{tab}{ide}

\olsection{\usetoken{P}{tableau} avec \usetoken{S}{identity}}

Les !!^{tableau}s avec !!{identity} exigent des règles d'inférence
supplémentaires. Les règles de~$\eq$ sont les suivantes, où $t$,
$t_1$ et $t_2$ sont des termes clos :

\begin{defish}
\AxiomC{}
\RightLabel{$\eq$}
\UnaryInfC{\sFmla{\True}{\eq[t][t]}}
\DisplayProof
\hfill
\AxiomC{\sFmla{\True}{\eq[t_1][t_2]}}
\noLine
\UnaryInfC{\sFmla{\True}{!A(t_1)}}
\RightLabel{$\TRule{\True}{\eq}$}
\UnaryInfC{\sFmla{\True}{!A(t_2)}}
\DisplayProof
\hfill
\AxiomC{\sFmla{\True}{\eq[t_1][t_2]}}
\noLine
\UnaryInfC{\sFmla{\False}{!A(t_1)}}
\RightLabel{$\TRule{\False}{\eq}$}
\UnaryInfC{\sFmla{\False}{!A(t_2)}}
\DisplayProof
\end{defish}
Contrairement à toutes les autres règles, $\TRule{\True}{\eq}$ et
$\TRule{\False}{\eq}$ exigent que \emph{deux} !!{signed formula}s
figurent déjà sur la branche, à savoir
$\sFmla{\True}{\eq[t_1][t_2]}$ et $\sFmla{S}{!A(t_1)}$.

\begin{ex}
Si $s$ et $t$ sont des termes clos, alors
$\eq[s][t], !A(s) \Proves !A(t)$ :
\begin{oltableau}
  [\sFmla{\False}{\formula{A}(t)}, just = \TAss
    [\sFmla{\True}{\eq[s][t]}, just = \TAss
      [\sFmla{\True}{\formula{A}(s)}, just = \TAss
        [\sFmla{\True}{\formula{A}(t)}, just={\TRule{\True}{\eq}[2, 3]}, close]
      ]
    ]
  ]
\end{oltableau}
On reconnaît ici le principe de substituabilité des identiques,
ou loi de Leibniz.

Les !!^{tableau}s montrent aussi que $\eq$ est symétrique, c'est-à-dire
que $\eq[s_1][s_2] \Proves \eq[s_2][s_1]$ :
\begin{oltableau}
  [\sFmla{\False}{\eq[s_2][s_1]}, just = \TAss
    [\sFmla{\True}{\eq[s_1][s_2]}, just = \TAss
      [\sFmla{\True}{\eq[s_1][s_1]}, just = {$\eq$}
        [\sFmla{\True}{\eq[s_2][s_1]}, just = {\TRule{\True}{\eq}[2, 3]}, close]
      ]
    ]
  ]
\end{oltableau}
Ici, la ligne~$2$ fournit la première !!{formula} requise par
$\TRule{\True}{\eq}$, à savoir
$\sFmla{\True}{\eq[s_1][s_2]}$. La ligne~$3$ fournit la seconde, de la
forme $\sFmla{\True}{!A(s_1)}$. Il suffit de prendre $!A(x)$ pour
$\eq[x][s_1]$ : $!A(s_1)$ est alors $\eq[s_1][s_1]$, et $!A(s_2)$
est $\eq[s_2][s_1]$.

Ils montrent également que $\eq$ est transitive, c'est-à-dire que
$\eq[s_1][s_2], \eq[s_2][s_3] \Proves \eq[s_1][s_3]$ :
\begin{oltableau}
  [\sFmla{\False}{\eq[s_1][s_3]}, just = \TAss
    [\sFmla{\True}{\eq[s_1][s_2]}, just = \TAss
      [\sFmla{\True}{\eq[s_2][s_3]}, just = \TAss
        [\sFmla{\True}{\eq[s_1][s_3]}, just = {\TRule{\True}{\eq}[3, 2]}, close]
      ]
    ]
  ]
\end{oltableau}
Dans ce !!{tableau}, la première !!{formula} requise par
$\TRule{\True}{\eq}$ est la ligne~$3$,
$\sFmla{\True}{\eq[s_2][s_3]}$ : $s_2$ joue le rôle de~$t_1$ et
$s_3$ celui de~$t_2$. La seconde, de la forme
$\sFmla{\True}{!A(s_2)}$, est la ligne~$2$. Prenons ici $!A(x)$ pour
$\eq[s_1][x]$ : $!A(s_2)$ devient $\eq[s_1][s_2]$, c'est-à-dire la
ligne~$2$, et $!A(s_3)$ devient la !!{formula}~$\eq[s_1][s_3]$ de la
conclusion.
\end{ex}

\begin{editorial}
Dans l'explication de la symétrie, la source donne par erreur la forme
$\sFmla{\True}{!A(s_2)}$ à la seconde formule requise de la ligne~$3$;
il s'agit de $\sFmla{\True}{!A(s_1)}$. Dans l'explication de la
transitivité, $!A(s_2)$ est rétabli comme $\eq[s_1][s_2]$, conformément
à la ligne~$2$, plutôt que comme l'égalité générique
$\eq[t_1][t_2]$ qui désigne ici la ligne~$3$.
\end{editorial}

\begin{prob}
Construire un !!{tableau} fermé pour chacun des ensembles suivants :
\begin{enumerate}
\item $\sFmla{\False}{\lforall[x][\lforall[y][((x = y \land !A(x))
      \lif !A(y))]]}$
\item $\sFmla{\False}{\lexists[x][(!A(x) \land
    \lforall[y][(!A(y) \lif y = x)])]}$,\\
  $\sFmla{\True}{\lexists[x][!A(x)] \land
    \lforall[y][\lforall[z][((!A(y) \land !A(z)) \lif y = z)]]}$
\end{enumerate}
\end{prob}

\end{document}
